%%\section{Introduction}

Accurate neutrino cross-section measurements and modeling of nuclear effects are required for precise measurements of neutrino oscillation physics such as CP-violation and the ordering of the neutrino masses~\cite{Abe:2011ks,Ayres:2004js,Acciarri:2015uup}. 
One of the most important channels for lower neutrino energy oscillation experiments is charged current quasielastic (CCQE) scattering, in which a neutrino exchanges a $W$ with a neutron, producing one lepton and one proton in the final state, along with a possibly excited nucleus which typically is undetected. Experiments such as T2K and MiniBooNE use CCQE interactions as the main channel for oscillation measurements ~\cite{Abe:2015ibe,Aguilar-Arevalo:2013pmq} because in principle the neutrino energy can be deduced using only the lepton kinematics and assuming a 2-body elastic scatter. Better understanding of the CCQE process as it occurs in a nuclear medium is also needed for the higher neutrino energies of NO$\nu$A~\cite{Adamson:2016tbq} and the future long baseline oscillation experiments DUNE and HyperK~\cite{Acciarri:2015uup,Ayres:2004js}.

Experiments used large, deuterium-filled bubble chambers in the 1970s--1990s to measure CCQE scattering on (quasi-) free nucleons and obtained consistent results \cite{Baker:1981su,Miller:1982qi,PhysRevD.28.436}. Recent experiments using heavier nuclei such as carbon, oxygen, and iron as targets have shown that our understanding of CCQE is incomplete~\cite{Gran:2006jn,AguilarArevalo:2008rc,AguilarArevalo:2010wv,PhysRevLett.111.022502,Fields:2013zhk,PhysRevD.91.012005}. Using heavier nuclei requires consideration of the role that the nuclear environment plays; the initial state neutrons are in a bound system, and the reaction products can interact on their way out of the nucleus.  A variety of initial state effects have been suggested, including Fermi motion~\cite{Smith:1972xh}, effects resulting in two particles and two holes (2p2h) such as meson exchange currents ~\cite{Martini:2009uj,Nieves:2013fr,PhysRevC.83.045501}, short range correlations ~\cite{Benhar:2006wy,Arrington:2011xs,Arrington:2012ax}, and long range correlations as estimated using the Random Phase Approximation (RPA) ~\cite{Graczyk:2003ru,Nieves:2004wx,Valverde:2006yi}. Final state interactions (FSI), where a hadron scatters on its way out of the nucleus, are modeled in a variety of implementations~\cite{Andreopoulos:2009rq,Nowak:2006xv,Golan:2013jtj}.

Nuclear effects modify both the final-state particle kinematics and content, altering the rate of detection of any given interaction channel.
The neutron's initial state affects the final-state particle kinematics, while FSI affect both kinematics and the particle content of the final state, since particles can be rescattered or absorbed in the nucleus. As a result, a sample of events with an observed lepton and nucleons in the final state may have originated via inelastic processes at production. For example,  $\Delta$ (1232)  resonance  production  and  decay,  if the pion is absorbed during FSI, leaves a QE-like final state that contains only one lepton and some number of nucleons, but no pions. In this case the neutrino energy reconstruction assuming simple two-body kinematics will be incorrect.  It is important that these $A$-dependent nuclear effects, which impact neutrino energy reconstruction as well as signal efficiencies, be understood and modeled since future oscillation experiments will use targets that range from carbon~\cite{Adamson:2016tbq} to argon~\cite{Acciarri:2015uup}.

 Previous measurements from \minerva used the kinematics of the muon or the proton to study quasielastic interactions on CH \cite{PhysRevLett.111.022502,Walton:2014esl}. Protons are affected by FSI, so measuring them provides new constraints on the FSI models. A direct way to elicit nuclear effects in neutrino interactions is by making simultaneous measurements on different nuclei in the same detector. This approach allows flux and detector uncertainties to be reduced by taking ratios of measurements on different nuclei. The first measurement of this kind to be based entirely upon CCQE-like interactions is presented here.

